\scanpage{120}
\deleted{Posons, pour $a\in A_1$,}
\[
\deleted{\varepsilon_a=\left\{\begin{gathered}\text{ce qui est défini pg. précédente, \unsure{si $s_0$ pas}}\\\text{un sommet de }a\end{gathered}\right.}
\]
et $J_1$ l'ens des éléments de $J=\widehat{\vec A}_{s_0}$ qui sont \deleted{\unsure{liés}} des arcs liés pas boucles, $J_2=J\setminus J_1$. On aura alors, posant
\[
\varepsilon_{\widehat{A}_1}=\prod_{a\in\widehat{A}_1}\varepsilon_a,
\]
\[
SP_X=\underbrace{\left(\mathbb{U}^{A_1}\wedge_{\{\pm1\}^{A_1}}\varepsilon_{A_1}\right)}_{T_{A_1}}\times T_J^Y,
\]
où
\[
T_J^Y=\prod_{a\in\widehat{A}\unsure{\setminus}\widehat{A}_1}T_a
\]
\margin{(produit d'au moins deux \deleted{\unsure{fbts}} au plus quatre facteurs}
\margin{$A\setminus A_1={}$ens des arêtes incidentes à $s_0$}

où $T_J^Y$ est \inserted{\unsure{ainsi}} défini directement en termes de $Y$ et des ens. \inserted{à deux éléments} $\widehat{\vec A}_s\setminus\{\vec a\}$ relatifs aux \deleted{sommets} \inserted{arcs $a\in\widehat{\vec A}_{s_0}=J$} qui sont liés et pas des boucles, $s$ désignant l'extrémité \deleted{\unsure{$s$ distinct de $s_0$ et adjacent à $s_0$}} dudit arc --- en \unsure{supposant} de plus \unsure{connu|connus} le groupage \unsure{d'|de}\inserted{certains} éléments de $J$ en \inserted{une ou deux} paires (disjointes dans le cas de deux paires), quand il y a une ou deux boucles de $G$ en $s_0$. De façon précise, par la structure de $G$ dont $J$ provient, on a
\[
J=J_{\mathrm{I}}\amalg J_{\mathrm{II}}\amalg J_{\mathrm{III}}
\]
où $J_{\mathrm{II}}$ est muni d'une involution sans pts fixes $\unsure{\sigma_{\mathrm{II}}}$.
et où $J_{\mathrm{III}}$ est muni d'un rev. de degré $2$
\diagram{\begin{tikzcd}\widetilde{J}_{\mathrm{III}}\arrow[r]&J_{\mathrm{III}}\end{tikzcd}}
On aura alors, \unsure{en termes de} cette structure de $J$
\[
T_J^Y=T_{J_{\mathrm{I}}}^Y\unsure{*|\times}T_{J_{\mathrm{II}}}^Y\unsure{*|\times}T_{J_{\mathrm{III}}}^Y
\]
\mspage{119}
